% Informe de Entrega 3 -- MIRA CU-1 (RAG + MCP)
% Estructura segun la pauta Entrega_3_MIRA_RAG_MCP.pdf, capitulo 11.
% Extensiones maximas: 11.1=0,5; 11.2=1,5; 11.3=2; 11.4=1,5; 11.5=1,5; 11.6=3; 11.7=1; 11.8=1 planas.
%
% ARCHIVOS: dejar en la MISMA carpeta que este .tex:
%   imagen1.png, imagen2.png y evidencia_01 ... evidencia_06 (.png)
%
%
\documentclass[11pt,a4paper]{article}
\usepackage[utf8]{inputenc}
\usepackage[T1]{fontenc}
\usepackage{lmodern}
\usepackage[spanish]{babel}
\usepackage[a4paper,left=1.9cm,right=1.9cm,top=1.8cm,bottom=1.9cm]{geometry}
\usepackage[table]{xcolor}
\usepackage{booktabs,tabularx,array}
\usepackage{enumitem}
\usepackage{titlesec}
\usepackage{graphicx}
\usepackage{float}
\usepackage[hidelinks]{hyperref}
\usepackage{microtype}
\usepackage{needspace}
\usepackage[font=small,labelfont=bf,skip=3pt]{caption}

\definecolor{azulmira}{RGB}{31,56,100}
\definecolor{turquesamira}{RGB}{13,124,120}
\definecolor{grisclaro}{RGB}{242,246,248}
\definecolor{rojoalerta}{RGB}{179,38,30}
\setlist{itemsep=1.5pt,topsep=3pt,parsep=0pt,leftmargin=1.3em}
\titlespacing*{\section}{0pt}{12pt}{5pt}
\titlespacing*{\subsection}{0pt}{8pt}{3pt}
\titleformat{\section}{\large\bfseries\color{azulmira}}{\thesection}{0.55em}{}[\color{azulmira}\titlerule]
\titleformat{\subsection}{\normalsize\bfseries\color{turquesamira}}{\thesubsection}{0.5em}{}
\setlength{\parskip}{4.5pt}
\setlength{\parindent}{0pt}
\setlength{\tabcolsep}{4pt}
\renewcommand{\arraystretch}{1.18}
\linespread{1.04}
\newcolumntype{Y}{>{\raggedright\arraybackslash}X}
\newcommand{\codigo}[1]{\texttt{\small #1}}
\newcommand{\completar}[1]{\textbf{\textcolor{rojoalerta}{[#1]}}}
\newcommand{\cab}{\rowcolor{azulmira!14}}
\newcommand{\tabla}{\small\rowcolors{2}{white}{grisclaro}}
\renewcommand{\contentsname}{Índice}
\setcounter{tocdepth}{1}

% Figura generica: si falta el archivo, deja un recuadro visible en lugar de fallar.
\newcommand{\figura}[4]{%
  \begin{figure}[#4]
  \centering
  \IfFileExists{#1.png}{\includegraphics[width=#3\linewidth,height=0.86\textheight,keepaspectratio]{#1.png}}{%
    \fbox{\parbox[c][3cm][c]{#3\linewidth}{\centering\textbf{Falta #1.png}}}}
  \caption{#2}
  \end{figure}}
\newcommand{\evidencia}[3]{%
  \begin{figure}[#3]
  \centering
  \IfFileExists{#1}{\includegraphics[width=0.9\linewidth,height=6.4cm,keepaspectratio]{#1}}{%
    \fbox{\parbox[c][2.7cm][c]{0.88\linewidth}{\centering\textbf{Falta \detokenize{#1}}}}}
  \caption{#2}
  \end{figure}}

\hypersetup{
  pdftitle={MIRA — Entrega 3: asistente con RAG y MCP},
  pdfauthor={Equipo MIRA},
  pdfsubject={Caso CU-1: asistente del operador de revisión}
}
\sloppy
\begin{document}
\pagestyle{plain}

% ---------------------------------------------------------------- PORTADA
\begin{titlepage}
\centering
{\large Universidad de Valparaíso}\\[2pt]
{\large Escuela de Ingeniería Informática}\\[2pt]
{\large Taller de Desarrollo de Software Conducido por Agentes de IA}\\[26pt]
{\Huge\bfseries\color{azulmira} MIRA}\\[7pt]
{\LARGE Entrega 3}\\[5pt]
{\Large Asistente con RAG y MCP:\\ evidencia documental y acceso gobernado}\\[6pt]
{\large Caso de uso CU-1: asistente del operador de revisión}\\[26pt]
\begin{tabularx}{0.9\textwidth}{@{}>{\bfseries}p{4.4cm}Y@{}}
Integrantes & Nikolas Lagos; Franko Moraga; Dante Chaves; Felipe Caru\\
Representante del cliente & Nikolas Lagos\\
Versión evaluada & v3.0\\
Fecha de entrega & 9 de noviembre de 2026\\
Repositorio & \url{https://github.com/Frxnks/Entregas_DSLLMs} (carpeta \codigo{entrega3})\\
Datos & Datos ficticios del kit CU-1\\
\end{tabularx}
\vfill
\fcolorbox{turquesamira}{grisclaro}{\parbox{0.84\textwidth}{\centering\vspace{2pt}
\textbf{Resultado principal}\\
La recuperación encuentra toda la evidencia esperada en 14 de 16 preguntas respondibles (87,5\,\%).\\
Las 20 preguntas del golden set se ejecutaron con Gemini sin errores técnicos y se revisaron citas y respuestas: 26 de las 27 fuentes nombradas respaldan lo afirmado y las 20 respuestas son fieles a la evidencia recibida. El manejo de conflictos y versiones es el punto débil (1 de 4).\vspace{2pt}}}
\vfill
\end{titlepage}

\tableofcontents
\clearpage

% ---------------------------------------------------------------- 11.1
\section{Resumen ejecutivo}
\textbf{Qué hace y para quién.} MIRA CU-1 es un asistente web para el operador de una organización cliente que revisa los casos que MIRA no resolvió sola. Le explica por qué un caso fue derivado, qué regla lo explica y qué dicen los documentos, y le prepara solicitudes de antecedentes. Nunca aprueba ni rechaza un caso: la decisión y la confirmación son de una persona.

\textbf{Qué tan bien funciona.} El asistente combina búsqueda documental (RAG) y un servidor propio de consultas (MCP). Al buscar en los documentos, encuentra toda la evidencia necesaria en 14 de 16 preguntas que tienen respuesta (87,5\,\%). Las 20 preguntas de prueba se ejecutaron con Gemini sin fallas técnicas; las 4 que no tienen respuesta en los documentos terminaron en abstención (4 de 4) y solo 1 de 16 preguntas respondibles se abstuvo sin motivo. En la revisión de las respuestas, 26 de las 27 fuentes que cada respuesta nombra respaldan lo afirmado (96\,\%) y las 20 respuestas son fieles a la evidencia que recibieron. El punto débil es el manejo de conflictos y versiones: solo 1 de 4 casos quedó bien resuelto.

\textbf{Qué no hace.} No decide casos, no responde cuando falta información (se abstiene) y, si dos documentos vigentes se contradicen, muestra ambos sin elegir. Falla cuando la búsqueda no trae las dos fuentes de un conflicto (C01) o de una comparación de versiones (V02), y en un caso (C02) trajo ambas pero respondió con un «No» categórico en vez de presentarlas sin elegir. No se declara corregido ninguno de los tres.

\textbf{Recomendación.} Usarlo como apoyo de lectura y preparación, no como autoridad. Antes de un uso real conviene corregir el conflicto de umbrales (80 vs.\ 85) en los documentos del cliente.

% ---------------------------------------------------------------- 11.2
\section{Caso de uso y especificación}
\subsection{Caso elegido y justificación}
Se implementa CU-1 (la pauta, sección 3.1). El operador pertenece a una organización ficticia, pregunta en español y necesita dos clases de respuesta que no pueden resolverse con una sola técnica: las \emph{reglas} están en documentos versionados (RAG) y los \emph{datos de un caso} viven en una base de datos (MCP). El kit se importa sin alterar sus filas: tres organizaciones, ocho usuarios, catorce casos, señales, evidencias, solicitudes y cambios de umbral.

\subsection{Contradicción del corpus}
El kit contiene una contradicción entre fuentes vigentes: POL-DEC v2 §2 fija 85 como umbral de referencia (sujeto a mínimos y configuración), mientras que MAN-REV v1 §1 indica aprobación desde 80 sin fraude. Ambas figuran vigentes; el asistente debe mostrarlas y no elegir (PR-4). POL-DEC v1 está derogada: solo sirve para explicar cambios históricos.

\subsection{Clasificación de las necesidades}
\begin{center}\tabla
\begin{tabularx}{\textwidth}{@{}>{\raggedright\arraybackslash}p{4.1cm}p{2.5cm}Y@{}}
\toprule
\cab\textbf{Necesidad del operador} & \textbf{Tipo} & \textbf{Argumento}\\
\midrule
¿Qué dice la política sobre el umbral y qué cambió entre versiones? & RAG & Conocimiento normativo en documentos versionados; exige cita y vigencia.\\
¿Cuál es el estado, las señales y las evidencias del caso C-012? & MCP (lectura) & Dato vivo de la base; el modelo no puede inventarlo (PR-5).\\
¿Por qué se derivó el caso C-012 con puntaje 86? & RAG + MCP & Regla del documento aplicada a los datos reales del caso.\\
Solicitar antecedentes sobre un caso & MCP (acción) & Efecto en la base: valida permisos y límites y exige confirmación (PR-6).\\
Preguntas que ningún documento responde (p.\,ej., costo de una solicitud) & Abstención & La ausencia de evidencia no se presenta como respuesta (PR-3).\\
\bottomrule
\end{tabularx}
\end{center}

\subsection{Flujos}
\textbf{Principal.} (1) El operador inicia sesión de prueba y pregunta. (2) La API valida la identidad por MCP y recupera reglas. (3) Si hacen falta datos del caso, se usan herramientas MCP de lectura. (4) La respuesta muestra estado, explicación, citas y traza. (5) Para una solicitud, MCP devuelve una propuesta; la persona la revisa y confirma; MCP revalida y registra.

\textbf{Alternativos.} \emph{Sin evidencia:} el asistente se abstiene y dice qué falta. \emph{Conflicto:} presenta ambas fuentes. \emph{Acción rechazada:} el servidor devuelve un error explicativo (rol sin permiso, caso aprobado automáticamente, tercera solicitud sin supervisor). \emph{Confirmación repetida:} es idempotente y no duplica. \emph{Proveedor o MCP caído:} la API informa el error (HTTP 503) sin inventar una respuesta.

\subsection{Requisitos trazables}
\begin{center}\tabla
\begin{tabularx}{\textwidth}{@{}p{1.7cm}Yp{1.5cm}p{2.9cm}@{}}
\toprule
\cab\textbf{ID} & \textbf{Requisito verificable} & \textbf{Principio} & \textbf{Fuente / prueba}\\
\midrule
RF-01 a 03 & Reglas con fuente, sección y versión; abstención si falta respaldo; conflicto presentado sin elegir. & PR-2, 3, 4 & POL-DEC, MAN-REV; golden set R, N, C\\
RF-04 a 07 & Datos y trazas por MCP; propuesta separada de la confirmación; control por organización y rol. & PR-5, 6, 7 & PRO-ANT §4; IT-01 a 06; D5 a D7\\
RNF-01 a 03 & Integridad del kit, sanitización de datos y errores explícitos. & PR-1, 7 & IT-13 a 15; suite Python\\
\bottomrule
\end{tabularx}
\end{center}
Los principios PR-1 a PR-7 de la pauta rigen el comportamiento del agente y se implementan en el servidor y la API, no solo en el prompt. El asistente no decide (PR-1).

% ---------------------------------------------------------------- 11.3
\section{Diseño e implementación del RAG}
\subsection{Corpus e inventario}
Se indexan los seis documentos obligatorios del kit. Los tokens son una estimación (caracteres UTF-8 divididos por cuatro), no el conteo del tokenizer de Gemini. No se agregaron documentos propios.

\begin{center}\tabla
\begin{tabularx}{\textwidth}{@{}p{1.9cm}Yp{1.5cm}p{1.7cm}p{1.6cm}p{1.5cm}@{}}
\toprule
\cab\textbf{Código} & \textbf{Título} & \textbf{Versión} & \textbf{Vigencia} & \textbf{Tokens} & \textbf{Origen}\\
\midrule
CAT-SEN & Catálogo de señales & 3 & Vigente & 780 & Kit\\
FAQ-OPE & Preguntas frecuentes & 2 & Vigente & 418 & Kit\\
MAN-REV & Manual de revisión & 1 & Vigente & 641 & Kit\\
POL-DEC & Política de decisión & 1 & Derogada & 381 & Kit\\
POL-DEC & Política de decisión & 2 & Vigente & 762 & Kit\\
PRO-ANT & Solicitud de antecedentes & 1 & Vigente & 521 & Kit\\
\bottomrule
\end{tabularx}
\end{center}
Cada fragmento conserva como metadatos código, título, versión, vigencia, estado, sección y líneas de origen. Los documentos derogados se indexan con su estado para poder explicar cambios, y el corpus se mantiene en Markdown para inspeccionarlo.

\subsection{Indexación y búsqueda}
La indexación y la consulta están separadas. \codigo{prepare.py} importa el Excel y construye el índice; \codigo{RagEngine.search} solo consulta. Se usa ChromaDB con distancia coseno y el modelo de embeddings \codigo{paraphrase-multilingual-MiniLM-L12-v2} (ONNX cuantizado). Se eligió por tres razones: el corpus y las preguntas están en español, las preguntas usan palabras distintas a las del documento (la búsqueda semántica las acerca), y corre en CPU sin costo ni envío del corpus a terceros. La búsqueda léxica existe solo para diagnóstico. La herramienta MCP \codigo{buscar\_\allowbreak reglas} expone la búsqueda; el parámetro de producción es \codigo{RAG\_TOP\_K=5}. El índice se reconstruye con un solo comando (anexo~\ref{anexo:repro}).

\subsection{Experimento de fragmentación}
Se compararon dos estrategias sobre las 16 preguntas respondibles del golden set (20 evidencias esperadas). Recall se mide a nivel de fragmento: la evidencia esperada debe aparecer entre los $k$ recuperados.

\begin{center}\tabla
\begin{tabularx}{\textwidth}{@{}p{3.1cm}p{1.6cm}p{2.3cm}p{2.3cm}Y@{}}
\toprule
\cab\textbf{Fragmentación} & \textbf{Fragm.} & \textbf{Recall@5 preg.} & \textbf{Recall@5 evid.} & \textbf{Decisión}\\
\midrule
Secciones completas & 35 & 14/16 = 0,875 & 18/20 = 0,90 & Elegida: mantiene juntas las condiciones de cada regla.\\
Ventanas de 450 caracteres, solape 90 & 40 & 14/16 = 0,875 & 18/20 = 0,90 & Sin ventaja demostrada; parte las reglas.\\
\bottomrule
\end{tabularx}
\end{center}
Con $k=4$ el Recall por pregunta era 0,8125 (13/16) en ambas estrategias; subir a $k=5$ recupera R05 y mantiene las dos evidencias de C02, sin ganar nada con ventanas. Se eligió secciones con $k=5$.

\subsection{Citas, abstención, conflicto y vigencia}
Cada cita muestra fuente, sección y versión, y la interfaz permite abrir el texto del fragmento. Si la evidencia recuperada no responde, el asistente se abstiene. Si recupera dos fuentes vigentes que se contradicen, presenta ambas y marca el estado de conflicto; si una versión está derogada, responde con la vigente e indica el cambio. La vigencia viaja en los metadatos y no se infiere solo por similitud. Limitación: un solo recuperador de $k$ fragmentos no siempre trae las dos fuentes de un conflicto (C01, sección~\ref{sec:fallos}).

% ---------------------------------------------------------------- 11.4
\section{Diseño e implementación del servidor MCP}
\subsection{Herramientas y contratos}
El servidor propio FastMCP opera por Streamable HTTP. Cada herramienta ejecuta una consulta parametrizada escrita por el equipo: el modelo elige la herramienta y sus argumentos y nunca escribe SQL. La base SQLite se carga desde el Excel del kit y solo la abre el servidor MCP.

\begin{center}\tabla
\begin{tabularx}{\textwidth}{@{}p{4.3cm}p{1.7cm}Y@{}}
\toprule
\cab\textbf{Herramienta} & \textbf{Tipo} & \textbf{Contrato (entrada $\rightarrow$ resultado)}\\
\midrule
\codigo{buscar\_\allowbreak reglas} & Lectura & Consulta de texto $\rightarrow$ fragmentos con fuente, sección, versión, vigencia y texto.\\
\codigo{consultar\_\allowbreak caso}, \codigo{listar\_\allowbreak senales}, \codigo{consultar\_\allowbreak evidencias} & Lectura & ID de caso $\rightarrow$ estado, señales o evidencias del caso, solo si pertenece a la organización del usuario.\\
\codigo{consultar\_\allowbreak solicitudes}, \codigo{consultar\_\allowbreak cambios\_\allowbreak umbral} & Lectura & ID de caso u organización $\rightarrow$ solicitudes previas o historial de umbrales.\\
\codigo{preparar\_\allowbreak solicitud\_\allowbreak antecedentes} & Propuesta & Caso, motivo, documentos $\rightarrow$ propuesta validada y token de un solo uso; \emph{no escribe}.\\
\codigo{solicitar\_\allowbreak antecedentes} & Escritura & Token de confirmación $\rightarrow$ revalida y registra; repetirlo no duplica.\\
\codigo{crear\_\allowbreak sesion}, \codigo{contexto\_\allowbreak sesion}, \codigo{consultar\_\allowbreak trazas} & Técnica & Sesión opaca, contexto del usuario y trazas sanitizadas.\\
\bottomrule
\end{tabularx}
\end{center}
Además hay un recurso con el inventario del corpus y el prompt \codigo{explicar\_\allowbreak caso}.

\subsection{Validaciones, errores e integridad}
La identidad se resuelve en una sesión opaca y la API inyecta el \codigo{session\_\allowbreak id}: el modelo no elige su organización ni su rol (PR-7). Antes de preparar o confirmar, el servidor valida sesión, organización, rol, estado del caso, documentos, motivo, límites y vigencia de la propuesta. Ejemplos: el usuario U-03 no puede registrar solicitudes; los casos aprobados automáticamente se bloquean; en C-016, que ya tiene dos solicitudes, la tercera exige supervisor (PRO-ANT §4) y la herramienta devuelve un mensaje que se muestra tal cual al usuario. La confirmación se revalida dentro de una transacción (\codigo{BEGIN IMMEDIATE}) y una consulta nunca produce efectos (PR-6).

La importación guarda una huella SHA-256 del Excel y de cada fila, y disparadores SQLite rechazan \codigo{UPDATE} y \codigo{DELETE} sobre los registros originales. Toda llamada deja una traza sanitizada con herramienta, argumentos, resultado y estado, visible en la interfaz (PR-5).

\subsection{Pruebas}
La matriz de integración IT-01 a IT-15 cubre roles, límites, aislamiento entre organizaciones, lecturas sin efectos, integridad, privacidad y fallas; la suite registra 28 pruebas aprobadas. Las pruebas con escritura usan una base aislada. Si MCP o el proveedor no responden, la API devuelve HTTP 503 con un mensaje claro.

% ---------------------------------------------------------------- 11.5
\section{Arquitectura e implementación front y back}
\subsection{Arquitectura}
El frontend Next.js (ruta \codigo{/asistente}) habla con una API FastAPI. La API valida la sesión por MCP, sanitiza el texto, coordina el proveedor Gemini y devuelve respuesta, citas y trazas. FastMCP autoriza cada llamada y consulta ChromaDB o SQLite según la herramienta. \textbf{La API no abre SQLite}; todo dato pasa por MCP. La secuencia completa, incluida la confirmación humana, está en el anexo~\ref{anexo:secuencia}.

\figura{imagen1}{Arquitectura técnica de MIRA CU-1. La API no accede directamente a SQLite; el acceso a los datos pasa por el servidor MCP.}{0.96}{htbp}

\subsection{Datos enviados al modelo}
\begin{center}\tabla
\begin{tabularx}{\textwidth}{@{}YY@{}}
\toprule
\cab\textbf{Se envía a Gemini} & \textbf{No se envía}\\
\midrule
Pregunta sanitizada (RUT, correo, teléfono y nombres conocidos enmascarados). Fragmentos del corpus recuperados, marcados como contenido no confiable. Resultados de herramientas MCP limitados a la consulta autorizada. & Claves, tokens de sesión y de confirmación, rutas del sistema, SQL y filas fuera de la organización del usuario.\\
\bottomrule
\end{tabularx}
\end{center}
La sanitización reduce la exposición accidental; no equivale a anonimizar.

\subsection{Interfaz y estado real del agente}
La interfaz en español indica de forma visible si la respuesta tiene evidencia, es una abstención, un conflicto o un error; permite abrir las citas y la traza de llamadas MCP, y una acción preparada se muestra completa antes de pedir «Confirmar y registrar» (anexo~\ref{anexo:capturas}). Si el proveedor o un servicio falla, se muestra un mensaje comprensible y no se inventa respuesta.

\textbf{Qué hace realmente el agente.} Gemini redacta con la evidencia recibida y puede pedir herramientas de lectura. La búsqueda de reglas no es opcional: la API la ejecuta por MCP antes de generar, porque en una primera prueba (R01) el modelo no la pidió. Así, la selección de esa herramienta la decide el software y no el modelo. La autorización, el aislamiento y la confirmación tampoco se delegan. El modo \codigo{demo} es determinista, se identifica en la interfaz y no se reporta como resultado de Gemini. El código se desarrolló con asistentes de IA, y el equipo revisó, probó y aprobó lo que entregó.

% ---------------------------------------------------------------- 11.6
\section{Evaluación}
\subsection{Estrategia}
Se separa recuperación de generación, de modo que cada fallo se pregunta primero si falló la búsqueda o falló el modelo. Hay cuatro niveles: (1) pruebas unitarias e integración de MCP, privacidad y fallas (28 pruebas); (2) experimento de recuperación sobre el golden set; (3) corrida de generación con Gemini (\codigo{gemini-flash-lite-latest}) sobre las 20 preguntas, cuyo reporte está en \codigo{evaluacion/generacion\_gemini\_20261108\_r01-r20.json}; (4) revisión de las 100 citas mostradas, de la fidelidad de las 20 respuestas y de los 4 casos de conflicto o versión, con la rúbrica binaria de la guía de revisión manual. Los juicios los hizo un modelo evaluador, con una nota por cada uno, y el equipo los verificó a mano sobre una muestra de 8 respuestas (C01, C02, V01, V02, R01, R05, N02 y N03), como exige la pauta cuando se usa un modelo evaluador; el equipo coincidió en los 8 casos.

\subsection{Golden set}
Lista fija de 20 preguntas con identificador, pregunta, fragmento esperado, respuesta esperada, si debe responder y tipo. Las respuestas esperadas las validó el representante del cliente (Nikolas Lagos).

\begin{center}\tabla
\begin{tabularx}{\textwidth}{@{}p{4.3cm}p{1.0cm}p{1.9cm}Y@{}}
\toprule
\cab\textbf{Tipo} & \textbf{N.º} & \textbf{IDs} & \textbf{Comportamiento esperado}\\
\midrule
Con respuesta en el corpus & 12 & R01--R12 & Responde y cita el fragmento esperado.\\
Sin respuesta en el corpus & 4 & N01--N04 & Se abstiene y dice qué falta (pregunta trampa).\\
Conflicto entre fuentes vigentes & 2 & C01--C02 & Presenta ambas y no elige.\\
Distinta respuesta entre versiones & 2 & V01--V02 & Responde con la versión vigente e indica el cambio.\\
\bottomrule
\end{tabularx}
\end{center}
\textbf{Limitación declarada.} La primera corrida de recuperación (7 de noviembre) es anterior a la revisión y validación del golden set del 8 de noviembre, en la que se corrigieron entradas para ajustarlas al corpus. Los resultados de recuperación y de generación aquí informados usan la versión validada, y no se modificó tras la validación. El historial queda en el repositorio.

\subsection{Resultados}
\begin{center}\tabla
\begin{tabularx}{\textwidth}{@{}p{3.3cm}p{3.6cm}Y@{}}
\toprule
\cab\textbf{Etapa / métrica} & \textbf{Resultado} & \textbf{Interpretación}\\
\midrule
Recall@5, por pregunta & 14/16 = 0,875 & Dos preguntas respondibles quedan con evidencia incompleta (C01 y V02).\\
Recall@5, por fragmento & 18/20 = 0,90 & Faltan 2 de 20 evidencias esperadas.\\
Ejecución con Gemini & 20/20; 0 errores & Las llamadas terminaron; no prueba que las respuestas sean correctas.\\
Abstención correcta & 4/4 (100\,\%) & Preguntas sin respuesta en que se abstuvo; resultado automático preliminar.\\
Abstención injustificada & 1/16 (6,25\,\%) & Solo V02, por evidencia insuficiente recuperada.\\
Precisión de cita, fragmentos nombrados & 26/27 (96,3\,\%) & Fuentes que la respuesta enumera; la única que no respalda es POL-DEC v1 (derogada) en C01.\\
Precisión de cita, todos los fragmentos mostrados & 28/100 (28\,\%) & La interfaz muestra los $k=5$ fragmentos recuperados aunque la respuesta use solo uno o dos; el resto es ruido de recuperación, no afirmaciones sin respaldo.\\
Fidelidad & 20/20 (100\,\%) & Todas las afirmaciones están respaldadas por la evidencia recibida (C01 es fiel pero incompleta porque no recibió MAN-REV).\\
Conflicto y vigencia & 1/4 (25\,\%) & Solo V01 correcta. C01 y V02 fallaron por recuperación; C02 expuso ambas fuentes pero concluyó.\\
Concordancia con la muestra humana & 8/8 (100\,\%) & Respuestas de la muestra en las que el equipo coincidió con el modelo evaluador.\\
\bottomrule
\end{tabularx}
\end{center}
Dos lecturas de la precisión de cita: la primera mide lo que la respuesta afirma haber usado; la segunda, todo lo que la interfaz muestra como referencia. La diferencia sugiere mostrar en la interfaz solo los fragmentos usados, o marcarlos, como mejora. La fidelidad de 100\,\% no equivale a corrección: una respuesta puede ser fiel a lo recibido y aun así estar incompleta.

\subsection{Clasificación de fallos}\label{sec:fallos}
\begin{center}\tabla
\begin{tabularx}{\textwidth}{@{}p{1.1cm}p{2.5cm}Yp{2.7cm}@{}}
\toprule
\cab\textbf{ID} & \textbf{Clasificación} & \textbf{Causa observada} & \textbf{Estado}\\
\midrule
R05 & Recuperación & Con $k=4$ no llegaba la evidencia esperada. & Resuelto con $k=5$\\
R01 & Recuperación (omisión) & En la primera corrida el modelo no consultó MCP (no pidió \codigo{buscar\_\allowbreak reglas}). & Resuelto: búsqueda obligatoria\\
C01 & Recuperación & No llegó MAN-REV v1 §1 (umbral 80); Gemini afirmó solo 85 y no expuso el conflicto. En modo demo, una búsqueda MCP adicional sí mostró ambas fuentes. & Sin mejora demostrada con Gemini\\
V02 & Recuperación & Gemini recibió POL-DEC v2 §5 y POL-DEC v1 §1, pero no la sección histórica necesaria para comparar quién podía modificar los umbrales; se abstuvo. & Sin corregir\\
C02 & Generación & Llegaron MAN-REV §1 y POL-DEC v2 §3, pero la respuesta concluyó «No» en lugar de presentar el conflicto sin elegir. & Sin corregir\\
\bottomrule
\end{tabularx}
\end{center}
Tras el ajuste, R01 mostró cinco referencias de las que dos respaldan la respuesta; la respuesta nombró solo esas dos, por lo que no se clasifica como fallo.

\subsection{Ciclo de mejora}
\textbf{Ciclo 1 (recuperación).} Hipótesis: $k=4$ deja fuera evidencia de reglas largas. Cambio: $k=5$. Antes: 13/16 (0,8125). Después: 14/16 (0,875); R05 se recupera y C02 mantiene sus dos evidencias.

\textbf{Ciclo 2 (generación y orquestación).} Hipótesis: pedir al modelo que use herramientas no garantiza que las use. Cambio: la API ejecuta \codigo{buscar\_\allowbreak reglas} antes de generar e incluye la evidencia en el contexto; se añadieron pruebas sobre búsqueda, citas y trazas. R01 se repitió y recuperó FAQ-OPE v2 con respuesta alineada.

\textbf{Ciclo 3 (C01, C02 y V02), sin mejora demostrada.} Con $k=5$, C01 y V02 recuperan una de sus dos evidencias esperadas. Subir $k$ de 4 a 5 sí corrigió R05, pero no C01 ni V02, y probar $k=6$ y $k=8$ tampoco mejoró el Recall. En C01, Gemini presentó POL-DEC v2 y no expuso el conflicto con MAN-REV v1, por lo que no cuenta como manejo correcto; en un ensayo aparte, el proveedor determinista hizo una búsqueda MCP adicional y mostró ambas fuentes, lo que sirve como evidencia del modo demo pero no corrige la evaluación de Gemini. En C02 la evidencia llegó completa y falló la redacción: la respuesta resolvió el conflicto en vez de exponerlo, por lo que corresponde ajustar la instrucción de conflicto del prompt y volver a medir. En V02 no se afirma mejora: se clasificó como fallo de recuperación. El siguiente paso es recuperar explícitamente la sección histórica de POL-DEC v1 y repetir la evaluación antes de declarar una mejora; el golden set no se modificará para ajustarlo al resultado.

% ---------------------------------------------------------------- 11.7
\section{Riesgos y controles}
\begin{center}\tabla
\begin{tabularx}{\textwidth}{@{}p{3.2cm}YY@{}}
\toprule
\cab\textbf{Riesgo} & \textbf{Control (guardrail)} & \textbf{Verificación}\\
\midrule
Inyección en un documento & Evidencia delimitada como contenido no confiable; permisos y decisiones no dependen del prompt. & Caso C-014 y pruebas, sin efectos de negocio.\\
Acceso cruzado o identidad falsa (PR-7) & Sesión opaca, identidad inyectada por la API, organización validada en MCP. & Pruebas U-01 y C-020.\\
Acción no autorizada o duplicada (PR-6) & Rol, estado y límites validados en el servidor; confirmación de un solo uso e idempotencia. & U-03, C-010, C-016, confirmación repetida.\\
Exposición de datos personales & Sanitización de RUT, correo, teléfono y nombres; exclusión de claves, tokens y rutas. & Pruebas de privacidad IT.\\
Cita incorrecta o evidencia faltante (PR-2, 3) & Citas expandibles, metadatos, abstención y revisión humana; los fallos C01 y V02 no se ocultan. & Golden set, guía de revisión manual y clasificación de fallos (sección~\ref{sec:fallos}).\\
Falla del proveedor o de MCP & Error explícito sin respuesta inventada; reintentos acotados. & Pruebas HTTP 503, IT-13.\\
Secretos y dependencias & La clave de Gemini solo va en el \codigo{.env} local; una clave apareció en un texto compartido y se recomendó revocarla y generar otra. \codigo{npm} informó 3 vulnerabilidades del frontend (1 moderada, 1 alta, 1 crítica), sin aplicar \codigo{audit fix --force}. & Guía de demostración y revisión de dependencias; la investigación de los avisos queda como mejora\\
\bottomrule
\end{tabularx}
\end{center}

% ---------------------------------------------------------------- 11.8
\section{Discusión y conclusiones}
\textbf{Qué resolvió cada componente.} RAG respondió lo normativo con cita y versión: por ejemplo, explicó el umbral y sus cambios, y se abstuvo ante preguntas sin respaldo. MCP entregó los datos vivos con permisos y trazas: U-03 no pudo registrar solicitudes y, en las pruebas, la tercera solicitud de C-016 exigió supervisor, con un mensaje que el asistente mostró tal cual. La confirmación humana cerró el ciclo: la propuesta no escribe y una repetición no duplica.

\textbf{Qué no resolvió ninguno.} El conflicto 80/85 entre POL-DEC v2 y MAN-REV v1 no se puede resolver con técnica: el sistema solo lo expone y corresponde al cliente decidir cuál rige (PR-4). Tampoco garantiza que la búsqueda traiga ambas fuentes: C01 y V02 mostraron que recuperar no equivale a recuperar lo suficiente: con Gemini el conflicto no se expuso (evidencia 4 del anexo) y solo se vio en modo demo. C02 mostró el fallo contrario: con la evidencia completa, el modelo igual eligió una de las dos fuentes.

\textbf{Aprendizajes transferibles.} (1) Una instrucción al modelo no garantiza una llamada de herramienta; lo crítico se fuerza en el software (caso R01). (2) Medir recuperación y generación por separado evita corregir el prompt cuando el fallo estaba en la búsqueda. (3) Aumentar $k$ ayuda poco si el problema es la composición del fragmento o la consulta. (4) Los límites y permisos se prueban en el servidor, no se piden en un prompt. (5) Un golden set se fija y se valida antes de medir; aquí la primera corrida de recuperación lo precedió, y lo declaramos.

\textbf{Pendientes concretos.} Que el cliente defina qué umbral rige y actualice el documento obsoleto; ampliar el golden set más allá de 20 preguntas; evaluar sistemáticamente la elección de herramientas con casos de uso real; ampliar la verificación humana de la revisión a las 20 respuestas; corregir la instrucción de conflicto para C02; recuperar la sección histórica de POL-DEC v1 para V02; y investigar las vulnerabilidades de npm. El recorrido D1--D8 se ensayó sobre la versión v3.0 y se presenta en vivo en la defensa.

\textbf{Uso de IA.} Se usó asistencia de IA para análisis, programación y documentación. El agente de la aplicación redacta y propone consultas, pero no decide ni confirma acciones. La selección del caso, la validación del golden set, la interpretación de los resultados y la aprobación final son del equipo.

\needspace{8\baselineskip}
\begin{thebibliography}{9}
\small
\addcontentsline{toc}{section}{Referencias}
\bibitem{pauta} Universidad de Valparaíso, \emph{Entregable 3: Asistente con RAG y MCP: evidencia documental y acceso gobernado al sistema}, pauta del caso MIRA, 2026.
\bibitem{kit} Kit CU-1, \emph{Operador de revisión}: corpus documental y datos ficticios entregados para la actividad.
\bibitem{repo} Equipo MIRA, \emph{Repositorio Entregas\_DSLLMs, carpeta entrega3}: código, golden set, resultados de recuperación y generación, matriz de integración y guía de demostración.
\end{thebibliography}

% ---------------------------------------------------------------- ANEXOS
\clearpage
\appendix
\section{Diagrama de secuencia}\label{anexo:secuencia}
\vspace{-6pt}
\figura{imagen2}{Secuencia de MIRA CU-1: pregunta documental y solicitud con confirmación humana. Gemini no autoriza ni escribe en la base; MCP controla el acceso a SQLite y la propuesta no se guarda antes de la confirmación.}{0.93}{H}

\clearpage
\section{Evidencias visuales}\label{anexo:capturas}
Las capturas usan datos ficticios y no contienen claves. La captura 4 documenta un primer intento con Gemini que no mostró el conflicto; no se considera una prueba exitosa. Las capturas 5 y 6 corresponden al modo demo, que sí presentó ambas fuentes con dos búsquedas MCP.

\evidencia{evidencia_01_interfaz_principal.png}{Interfaz inicial con el usuario ficticio U-01.}{H}
\evidencia{evidencia_02_respuesta_con_evidencia.png}{Respuesta con cita documental y traza MCP desplegada.}{H}
\evidencia{evidencia_03_abstencion_costo_antecedentes.png}{Abstención ante una consulta cuyo costo no está informado en el corpus.}{H}
\evidencia{evidencia_04_consulta_umbral_versiones.png}{Primer intento con Gemini del umbral: solo aparece POL-DEC v2 y el conflicto no queda demostrado.}{H}
\evidencia{evidencia_05_conflicto_umbral_demo.png}{Consulta del umbral en modo demo: estado CONFLICTO y explicación neutral de las dos fuentes.}{H}
\evidencia{evidencia_06_referencias_y_traza.png}{Referencias desplegadas y dos llamadas \codigo{buscar\_\allowbreak reglas} en la traza (modo demo).}{H}

\section{Reproducibilidad}\label{anexo:repro}
Dentro de \codigo{entrega3/service}: (1) crear el entorno con Python 3.12 e instalar el paquete; (2) copiar \codigo{.env.example} a \codigo{.env} (con \codigo{LLM\_PROVIDER=demo} si no hay clave de Gemini); (3) \codigo{python -m mira\_cu1.prepare} importa el Excel e indexa el corpus; (4) \codigo{python -m mira\_cu1.mcp\_server} y \codigo{python -m mira\_cu1.main} levantan MCP y API; (5) iniciar el frontend como indica \codigo{GUIA\_DEMO.md}. Docker no es necesario.

\end{document}
